\documentclass[../../../include/open-logic-section]{subfiles}

\begin{document}

\olfileid{sth}{card-arithmetic}{opps}
\olsection{মৌলিক ক্রিয়াগুলির সংজ্ঞা}

এখানে ক্রমের হিসাব রাখতে হয় না বলে অঙ্কবাচক সংখ্যার
পাটিগণিতের সংজ্ঞা ক্রমসংখ্যার পাটিগণিতের তুলনায় সহজ।
যোগ, গুণ ও ঘাত একসঙ্গে সংজ্ঞায়িত করব।

\begin{defn}
$\cardfont{a}$ ও $\cardfont{b}$ অঙ্কবাচক সংখ্যা হলে:
\begin{align*}
	\cardfont{a} \cardplus \cardfont{b} &\defis
	\card{\cardfont{a} \disjointsum \cardfont{b}}\\
	\cardfont{a} \cardtimes \cardfont{b} &\defis
	\card{\cardfont{a} \times \cardfont{b}}\\
	\cardexpo{\cardfont{a}}{\cardfont{b}} &\defis
	\card{\funfromto{\cardfont{b}}{\cardfont{a}}}
\end{align*}
এখানে $\funfromto{X}{Y} = \Setabs{f}{f \text{ is a function } X \to Y}$।
(যেকোনো সেট $X$ ও $Y$-এর জন্য $\funfromto{X}{Y}$ আছে,
এটি সহজেই দেখানো যায়; অনুশীলন হিসেবে রইল।)
\end{defn}

\begin{prob}
যেকোনো সেট $X$ ও $Y$-এর জন্য $\funfromto{X}{Y}$ আছে,
এটি $\Zminus$-এ প্রমাণ করো। $\ZF$-এ কাজ করে
$\setrank{X}$ ও $\setrank{Y}$ থেকে
$\setrank{\funfromto{X}{Y}}$ নির্ণয় করো,
\olref[sth][ord-arithmetic][using-addition]{rankcomputation}-এর মতো।
\end{prob}

সংজ্ঞাটি একটু ব্যাখ্যা করা যাক। যোগের ক্ষেত্রে এখানে
\olref[ord-arithmetic][add]{defdissum}-এ সংজ্ঞায়িত
বিযুক্ত যোগফল $\disjointsum$ ব্যবহার করেছি; সসীম
ক্ষেত্রে এর ফল যে ঠিক, তা সহজেই দেখা যায়। গুণের
ক্ষেত্রে \olref[sfr][set][pai]{cardnmprod} বলে,
$A$-র $n$ জন এবং $B$-র $m$ জন সদস্য থাকলে
$A\times B$-র $n\cdot m$ জন সদস্য। আমাদের
সংজ্ঞা এই ধারণাকেই ট্রান্সফাইনাইট গুণে প্রসারিত করে।
ঘাতও একইভাবে সসীম ধারণাটিকে ট্রান্সফাইনাইট ক্ষেত্রে
প্রসারিত করে। আসলে ট্রান্সফাইনাইট অঙ্কবাচক সংখ্যার
পাটিগণিত কিছু দিক থেকে ট্রান্সফাইনাইট ক্রমসংখ্যার
পাটিগণিতের চেয়ে ``সাধারণ'' পাটিগণিতের বেশি কাছাকাছি:

\begin{prop}\ollabel{cardplustimescommute}
$\cardplus$ ও $\cardtimes$ বিনিময়ী এবং সম্বন্ধীয়।
\end{prop}

\begin{proof}
বিনিময়িতার জন্য
\olref[cardinals][cardsasords]{lem:CardinalsBehaveRight}
অনুযায়ী $\cardeq{(\cardfont{a} \disjointsum
\cardfont{b})}{(\cardfont{b} \disjointsum \cardfont{a})}$
এবং $\cardeq{(\cardfont{a} \times \cardfont{b})}{(\cardfont{b} \times
\cardfont{a})}$ লক্ষ করলেই চলে। সম্বন্ধীয়তা অনুশীলন হিসেবে রইল।
\end{proof}

\begin{prob}
$\cardplus$ ও $\cardtimes$ সম্বন্ধীয়—প্রমাণ করো।
\end{prob}

\begin{prop}
$A$ অসীম তখনই, যখন $\card{A} \cardplus 1 = 1 \cardplus \card{A} = \card{A}$।
\end{prop}

\begin{proof}
\olref[cardinals][classing]{generalinfinitycharacter}-এর মতো,
\olref[ord-arithmetic][using-addition]{ordinfinitycharacter}
ও \olref[cardinals][cardsasords]{lem:CardinalsBehaveRight} থেকে।
\end{proof}

এতে বোঝা যায়, ক্রমসংখ্যা ও অঙ্কবাচক সংখ্যার যোগ ও
গুণের জন্য আলাদা চিহ্ন কেন দরকার: এগুলি সত্যিই
\emph{ভিন্ন} ক্রিয়া। পরের দুটি ফল দেখাবে, তাদের
ঘাতও ভিন্ন ক্রিয়া। (মনে করো,
\olref[z][infinity-again]{defnomega} থেকে $2=\{0,1\}$):

\begin{lem}\ollabel{lem:SizePowerset2Exp}
যেকোনো $A$-র জন্য $\card{\Pow{A}} = \cardexpo{2}{\card{A}}$।
\end{lem}

\begin{proof}
প্রত্যেক উপসেট $B\subseteq A$-র জন্য
$\chi_B\in\funfromto{A}{2}$ এভাবে নিই:
\begin{align*}
	\chi_{B}(x) &\defis
	\begin{cases}
		1 & \text{if }x\in B\\
		0 & \text{otherwise.}
	\end{cases}
\end{align*}
এখন $f(B)=\chi_B$ নিলে একটি !!a{bijection}
$f\colon\Pow{A}\to\funfromto{A}{2}$ পাই। তাই
$\cardeq{\Pow{A}}{\funfromto{A}{2}}$। এর ফলে
$\cardeq{\Pow{A}}{\funfromto{\card{A}}{2}}$;
সুতরাং $\card{\Pow{A}} = \card{\funfromto{\card{A}}{2}} =
2^{\card{A}}$।
\end{proof}

এই সংক্ষিপ্ত প্রমাণটি কার্যত
\olref[sfr][siz][red-alt]{sec}-এর আলোচনা ধারণ করে।
সেখানে দেখিয়েছিলাম কীভাবে $\Pow{\omega}$-র
অগণনীয়তার প্রশ্নটি অসীম বাইনারি অনুক্রমের সেট
$\Bin^\omega$-র অগণনীয়তার প্রশ্নে নামিয়ে আনা যায়।
কার্যত $\Bin^{\omega}$ হলো $\funfromto{\omega}{2}$;
উপরের প্রমাণ দেখায়, \olref[sfr][siz][red-alt]{sec}-এর
যুক্তি $\omega$-র বদলে যেকোনো সেট $A$-র জন্যই চলে।
এই ফল থেকে অঙ্কবাচক ঘাত সম্বন্ধেও একটি সহজ কথা পাই:

\begin{cor}\ollabel{cantorcor}
যেকোনো অঙ্কবাচক সংখ্যা $\cardfont{a}$-র জন্য
$\cardfont{a} < \cardexpo{2}{\cardfont{a}}$।
\end{cor}

\begin{proof}
কান্টরের উপপাদ্য (\olref[sfr][siz][car]{thm:cantor}) এবং
\olref{lem:SizePowerset2Exp} থেকে।
\end{proof}
\noindent
তাই $\omega<\cardexpo{2}{\omega}$। তবে এটি
\emph{অঙ্কবাচক সংখ্যার} ঘাতের ফল। \emph{ক্রমসংখ্যার}
ঘাতের সঙ্গে তুলনা করো: সেখানে
$\omega=\ordexpo{2}{\omega}$
(দেখো \olref[ord-arithmetic][expo]{sec})।

অঙ্কবাচক সংখ্যার ঘাতের প্রসঙ্গেই $\Real$ যে
!!{nonenumerable}, তা আরেকটু নির্দিষ্টভাবে বলা যায়।

\begin{thm}\ollabel{continuumis2aleph0}
$\card{\Real}=\cardexpo{2}{\omega}$।
\end{thm}

\begin{proof}[প্রমাণের রূপরেখা]
এটি প্রমাণের অনেক উপায় আছে। সহজ উপায় হলো
$\cardle{\Pow{\omega}}{\Real}$ এবং
$\cardle{\Real}{\Pow{\omega}}$ দেখানো। তারপর
শ্র্যোডার--বার্নস্টাইন থেকে
$\cardeq{\Real}{\Pow{\omega}}$ এবং
\olref[card-arithmetic][opps]{lem:SizePowerset2Exp}
থেকে $\card{\Real}=\cardexpo{2}{\omega}$ পাওয়া যায়।
$f\colon\Pow{\omega}\to\Real$ ও
$g\colon\Real\to\Pow{\omega}$ অন্তঃস্থাপন দুটি
সংজ্ঞায়িত করা একটি শিক্ষণীয় অনুশীলন হিসেবে রইল।
\end{proof}

\begin{prob}
$\cardle{\Pow{\omega}}{\Real}$ এবং
$\cardle{\Real}{\Pow{\omega}}$ দেখিয়ে
\olref[sth][card-arithmetic][opps]{continuumis2aleph0}-এর
প্রমাণ সম্পূর্ণ করো।
\end{prob}

\end{document}
